%% kapa: exam=1 round=33 year=2022 period=1 form=A subject=real_estate subject_ko=부동산학원론 seq=35 number=115
%% source: Q-Net 감정평가사 기출문제 (cst00309) artlSeq=5209997 file=raw/1008/5209997/2022년도 제33회 감정평가사 1차시험 1교시(민법,경제학,부동산학).pdf
%% conversion: kordoc 4.24.51 -> tools/gen1.py (auto); review=proofread-sonnet; flags=none
\begin{question}{115}
감정평가에 관한 규칙상 원가방식에 관한 설명으로 옳지 않은 것은?
\begin{choices}
  \item 원가법과 적산법은 원가방식에 속한다.
  \item 적산법에 의한 임대료 평가에서는 대상물건의 재조달원가에 기대이율을 곱하여 산정된 기대수익에 대상물건을 계속하여 임대하는 데에 필요한 경비를 더한다.
  \item 원가방식을 적용한 감정평가서에는 부득이한 경우를 제외하고는 재조달원가 산정 및 감가수정 등의 내용이 포함되어야 한다.
  \item 입목 평가 시 소경목림(小徑木林)인 경우에는 원가법을 적용할 수 있다.
  \item 선박 평가 시 본래 용도의 효용가치가 있으면 선체·기관·의장(艤裝)별로 구분한 후 각각 원가법을 적용해야 한다.
\end{choices}
\end{question}
\answer{2}
